\section{Extraction of Individual Diffusion Coefficients}
\label{sec:extraction}

Having established the composite timelag, steady-state flux, and upstream lead
time expressions, we now address the inverse problem: given experimental
measurements on a bilayer membrane, how can the individual transport
parameters be determined?

\subsection{The Inverse Problem}
\label{sec:inverse-problem}

For a bilayer membrane, four material parameters are unknown a priori:
\begin{equation}
  \label{eq:unknowns}
  \Ds{1},\quad \Ds{2},\quad \Ss{1},\quad \Ss{2}.
\end{equation}
The layer thicknesses $\Ls{1}$ and $\Ls{2}$ are assumed known from fabrication
records or cross-sectional microscopy.  The applied upstream pressure~$P$ is a
controlled experimental quantity.

The key structural observation is that the observables separate into two
groups:
\begin{enumerate}
  \item \textbf{Time-domain quantities} ($\tlag$, $\theta_{\mathrm{up}}$)
        depend on $\Ds{1}$, $\Ds{2}$, and the partition coefficient
        $\Kpart = \Ss{2}/\Ss{1}$, but \emph{not} on $\Ss{1}$ or $\Ss{2}$
        individually (\cref{eq:tlag-composite,eq:theta-up-explicit}).
  \item \textbf{Flux quantities} ($\Jss$) depend on the permeabilities
        $\Phi_i = \Ds{i}\Ss{i}$ (\cref{eq:Jss-series}), coupling diffusivity
        and solubility.
\end{enumerate}

\begin{remark}[Direction independence]
By \cref{prop:direction-independence}, the downstream timelag $\tlag$ is the
same in both directions.  The steady-state flux $\Jss$ is likewise
direction-independent (\cref{eq:Jss-series}).  Consequently, reversing the
sample orientation yields no new information from $\tlag$ or $\Jss$.
The upstream lead time $\theta_{\mathrm{up}}$ \emph{is}
direction-dependent (\cref{prop:upstream-direction}), but still depends only on
$\Ds{1}$, $\Ds{2}$, $\Kpart$.
\end{remark}

\paragraph{Counting equations and unknowns.}
From a single permeation experiment, we can extract at most three independent
quantities: $\tlag$, $\Jss$, and (if the upstream transient is measured)
$\theta_{\mathrm{up}}$.  These provide:
\begin{itemize}
  \item $\tlag = f(\Ds{1}, \Ds{2}, \Kpart)$ \hfill (1~equation in 3~unknowns)
  \item $\theta_{\mathrm{up}} = g(\Ds{1}, \Ds{2}, \Kpart)$
        \hfill (1~equation in same 3~unknowns)
  \item $\Jss = h(\Phi_1, \Phi_2) = h(\Ds{1}\Ss{1}, \Ds{2}\Ss{2})$
        \hfill (1~equation in 2~additional unknowns)
\end{itemize}
With four unknowns ($\Ds{1}$, $\Ds{2}$, $\Ss{1}$, $\Ss{2}$) and at most three
equations from one experiment, the system is \emph{under-determined}.
Additional information is required, and we describe several strategies below.

\subsection{Method~1: Single-Layer Calibration + Bilayer Experiment}
\label{sec:method1}

The most direct approach is to independently characterise one of the two
materials in a separate single-layer permeation experiment.

\paragraph{Step~1: Single-layer experiment on material~$i$.}
A single-layer sample of material~$i$ (thickness~$\Ls{i}^{*}$, where the
asterisk denotes the calibration sample) yields:
\begin{align}
  \label{eq:single-cal-D}
  \Ds{i} &= \frac{(\Ls{i}^{*})^2}{6\,\tau_i^{*}},\\[4pt]
  \label{eq:single-cal-Phi}
  \Phi_i &= \frac{\Jss^{*}\,\Ls{i}^{*}}{\sqrt{P}},
  \qquad\text{whence}\qquad
  \Ss{i} = \frac{\Phi_i}{\Ds{i}}.
\end{align}
This determines $\Ds{i}$ and $\Ss{i}$ for one material.

\paragraph{Step~2: Bilayer experiment.}
With $\Ds{i}$ and $\Ss{i}$ known, the bilayer timelag and flux provide two
equations in the two remaining unknowns $\Ds{j}$ and $\Ss{j}$ ($j \neq i$).

From \cref{eq:Jss-series}:
\begin{equation}
  \label{eq:method1-flux}
  \Phi_j = \Ds{j}\Ss{j}
  = \frac{\Ls{2}}{\dfrac{\sqrt{P}}{\Jss} - \dfrac{\Ls{1}}{\Phi_1}}.
\end{equation}
(Here we assumed $i=1$; the formula for $i=2$ is analogous.)

This determines $\Phi_j$.  To separate $\Ds{j}$ and $\Ss{j}$, we use the
timelag.  With $\Kpart = \Ss{2}/\Ss{1}$ and $\Phi_j = \Ds{j}\Ss{j}$, the
composite timelag \cref{eq:tlag-composite} becomes a single equation in
$\Ds{j}$ (since $\Kpart = \Phi_j/(\Ds{j}\Ss{i}) \cdot (\Ss{i}/\Ss{j})$
reduces to known quantities once $\Ds{j}$ is fixed).

More explicitly, define $\Kpart = \Ss{j}\Ss{i}^{-1} = \Phi_j/(\Ds{j}\Ss{i})$
(for $j=2$, $i=1$).  Substituting into \cref{eq:tlag-composite} gives a
single nonlinear equation in the one unknown $\Ds{j}$, which can be solved
numerically.  Once $\Ds{j}$ is found, $\Ss{j} = \Phi_j/\Ds{j}$.

\begin{proposition}[Solvability of Method~1]
\label{prop:method1-solvability}
Given the single-layer calibration of material~$i$ and the bilayer measurements
$\tlag$ and $\Jss$, the remaining parameters $\Ds{j}$ and $\Ss{j}$ are
generically determined uniquely (subject to the physical constraint
$\Ds{j} > 0$, $\Ss{j} > 0$).
\end{proposition}

\begin{proof}[Proof sketch]
The flux equation fixes $\Phi_j > 0$.  With $\Phi_j$ known, $\Kpart =
\Phi_j/(\Ds{j}\Ss{i})$ is a decreasing function of $\Ds{j}$.  Substituting
into the timelag formula yields a continuous function $\tlag(\Ds{j})$ on
$(0,\infty)$.  Since $\tlag \to \infty$ as $\Ds{j} \to 0$ and $\tlag \to
\Ls{i}^2/(6\Ds{i})$ as $\Ds{j} \to \infty$ (the fast-layer limit), the
intermediate value theorem guarantees at least one solution.  Monotonicity of
$\tlag(\Ds{j})$ on the feasible domain (which can be verified by examining the
derivative) ensures uniqueness.
\end{proof}

\begin{remark}
Method~1 requires a separate single-layer sample of one material.  This is
often available in practice: for hydrogen permeation barriers, the
\emph{substrate} material (e.g., a steel) is typically well-characterised,
and only the \emph{coating} or \emph{liner} is novel.
\end{remark}

\subsection{Method~2: Two Bilayer Samples with Different Thickness Ratios}
\label{sec:method2}

When single-layer calibration is not feasible (e.g., the individual materials
cannot be fabricated as free-standing membranes), one can fabricate two bilayer
samples of the \emph{same material pair} but with different thickness ratios.

\paragraph{Setup.}
Let sample~$k$ ($k = 1,2$) have thicknesses $\Ls{1}^{(k)}$ and
$\Ls{2}^{(k)}$, with $\Ls{1}^{(1)}/\Ls{2}^{(1)} \neq
\Ls{1}^{(2)}/\Ls{2}^{(2)}$.  From each sample, measure $\tlag^{(k)}$ and
$\Jss^{(k)}$.  This gives four equations:
\begin{align}
  \label{eq:method2-tlag}
  \tlag^{(k)} &= f\!\left(\Ds{1},\Ds{2},\Kpart;\,
    \Ls{1}^{(k)},\Ls{2}^{(k)}\right),
    \quad k=1,2,\\[4pt]
  \label{eq:method2-Jss}
  \Jss^{(k)} &= \frac{\sqrt{P}}{\dfrac{\Ls{1}^{(k)}}{\Phi_1}
    + \dfrac{\Ls{2}^{(k)}}{\Phi_2}},
    \quad k=1,2.
\end{align}

\paragraph{Reduction.}
The two flux equations \cref{eq:method2-Jss} are linear in the resistivities
$1/\Phi_i$:
\begin{equation}
  \label{eq:method2-flux-system}
  \frac{\Ls{1}^{(k)}}{\Phi_1} + \frac{\Ls{2}^{(k)}}{\Phi_2}
  = \frac{\sqrt{P}}{\Jss^{(k)}}, \quad k=1,2.
\end{equation}
Since the coefficient matrix
$\bigl(\begin{smallmatrix}\Ls{1}^{(1)}&\Ls{2}^{(1)}\\
\Ls{1}^{(2)}&\Ls{2}^{(2)}\end{smallmatrix}\bigr)$ is nonsingular whenever the
thickness ratios differ, this $2\times 2$ linear system determines $\Phi_1$ and
$\Phi_2$ uniquely.

With $\Phi_1$, $\Phi_2$ known, the two timelag equations
\cref{eq:method2-tlag} become two equations in three unknowns ($\Ds{1}$,
$\Ds{2}$, $\Kpart$).  However, $\Kpart = \Ss{2}/\Ss{1} =
(\Phi_2/\Ds{2})/(\Phi_1/\Ds{1}) = \Phi_2\Ds{1}/(\Phi_1\Ds{2})$, reducing the
system to two equations in two unknowns ($\Ds{1}$, $\Ds{2}$), which is
generically solvable.  The solubilities then follow:
$\Ss{i} = \Phi_i/\Ds{i}$.

\begin{proposition}[Solvability of Method~2]
\label{prop:method2-solvability}
Two bilayer samples with distinct thickness ratios and known layer thicknesses
generically determine all four parameters $\Ds{1}$, $\Ds{2}$, $\Ss{1}$,
$\Ss{2}$ from the measured timelags and steady-state fluxes.
\end{proposition}

\begin{remark}
With $N > 2$ samples, the system is over-determined.  A least-squares approach
\begin{equation}
  \label{eq:method2-lsq}
  \min_{\Ds{1},\Ds{2},\Ss{1},\Ss{2}>0}\;\sum_{k=1}^{N}
  \left[w_\tau\!\left(\tau^{(k)}_{\mathrm{meas}}
    - \tau^{(k)}_{\mathrm{model}}\right)^{\!2}
  + w_J\!\left(J^{(k)}_{\mathrm{ss,meas}}
    - J^{(k)}_{\mathrm{ss,model}}\right)^{\!2}\right]
\end{equation}
(with suitable weights $w_\tau$, $w_J$) provides uncertainty estimates and can
detect systematic discrepancies indicating, e.g., interface resistance or
concentration-dependent diffusivity.
\end{remark}

\subsection{Method~3: Full Transient Curve Fitting}
\label{sec:method3}

The cumulative permeation curve $\Cum(t)$ contains more information than the
three asymptotic parameters ($\tlag$, $\Jss$, $\theta_{\mathrm{up}}$).  In
particular, the \emph{shape} of the transient region encodes the eigenvalue
structure of the composite diffusion problem, which depends on $\Ds{1}$,
$\Ds{2}$, $\Kpart$, $\Ls{1}$, and $\Ls{2}$.

\paragraph{Normalised transient.}
Define the normalised cumulative permeation
\begin{equation}
  \label{eq:normalised-Q}
  \hat{\Cum}(t) \coloneqq \frac{\Cum(t)}{\Jss\,t},
\end{equation}
which tends to~$1$ at large~$t$ and whose time evolution depends only on
$\Ds{1}$, $\Ds{2}$, $\Kpart$ (and the known thicknesses), \emph{not} on the
absolute solubilities.

\paragraph{Four-parameter fit.}
Fitting the full (unnormalised) $\Cum(t)$ to the analytical model
(\cref{eq:Qbar-two-layer} inverted numerically, or equivalently a numerical PDE
solution) determines all four parameters simultaneously:
\begin{itemize}
  \item The transient shape pins down $\Ds{1}$, $\Ds{2}$, and $\Kpart$.
  \item The absolute magnitude of $\Jss$ then determines $\Ss{1}$ (and hence
        $\Ss{2} = \Kpart\,\Ss{1}$) via \cref{eq:Jss-series}.
\end{itemize}

\begin{remark}
This method extracts all four parameters from a \emph{single} experiment but
requires high-quality time-resolved data throughout the transient.  It is also
more computationally demanding and may be sensitive to noise, particularly when
the two layers have similar properties.  Regularisation or Bayesian inference
approaches can improve robustness.  Numerical validation with FESTIM is
recommended to verify the extracted parameters.
\end{remark}

\subsection{Special Case: Known Solubilities}
\label{sec:known-solubilities}

When the solubility coefficients $\Ss{1}$ and $\Ss{2}$ are known from
independent measurements (literature values, separate sorption experiments, or
electrochemical methods), the problem reduces significantly: only $\Ds{1}$ and
$\Ds{2}$ are unknown, and a single bilayer experiment suffices.

\paragraph{Method~A: Timelag + steady-state flux.}
From \cref{eq:Jss-series}, with $\Phi_i = \Ds{i}\Ss{i}$:
\begin{equation}
  \label{eq:extract-Jss-linear}
  \frac{\Ls{1}}{\Ds{1}\Ss{1}} + \frac{\Ls{2}}{\Ds{2}\Ss{2}}
  = \frac{\sqrt{P}}{\Jss}.
\end{equation}
This is one linear equation in $1/\Ds{1}$ and $1/\Ds{2}$.  The composite
timelag \cref{eq:tlag-composite} provides a second (nonlinear) equation.

Define $u_i \coloneqq 1/\Ds{i}$.  From \cref{eq:extract-Jss-linear}:
\begin{equation}
  \label{eq:u2-from-u1}
  u_2 = \frac{\Ss{2}}{\Ls{2}}
  \left(\frac{\sqrt{P}}{\Jss} - \frac{\Ls{1}}{\Ss{1}}\,u_1\right).
\end{equation}
Substituting into \cref{eq:tlag-composite} yields a single equation in $u_1$,
solvable numerically.

\paragraph{Method~B: Timelag + upstream lead time.}
If $\theta_{\mathrm{up}}$ is measurable, the pair $(\tlag,
\theta_{\mathrm{up}})$ provides two equations in $(\Ds{1}, \Ds{2})$ with
$\Kpart = \Ss{2}/\Ss{1}$ known.  This has the advantage of relying only on
\emph{time-domain} intercepts, avoiding the need for absolute flux calibration.

\subsection{Dimensionless Diagnostics and Conditioning}
\label{sec:conditioning}

Rather than examining the six individual parameters ($\Ds{1}$, $\Ds{2}$,
$\Ss{1}$, $\Ss{2}$, $\Ls{1}$, $\Ls{2}$) separately, the conditioning of the
extraction is governed by two dimensionless ratios that capture the
\emph{relative} contributions of the two layers.

\subsubsection{The Governing Dimensionless Groups}

\begin{definition}[Resistance ratio]
\label{def:kappa}
The \emph{resistance ratio} is
\begin{equation}
  \label{eq:kappa-def}
  \kres \coloneqq \frac{R_1}{R_2}
  = \frac{\Ls{1}/\Phi_1}{\Ls{2}/\Phi_2}
  = \frac{\Ls{1}\,\Ds{2}\,\Ss{2}}{\Ls{2}\,\Ds{1}\,\Ss{1}},
\end{equation}
where $R_i = \Ls{i}/\Phi_i = \Ls{i}/(\Ds{i}\Ss{i})$ is the permeation
resistance of layer~$i$.
\end{definition}

$\kres$ determines \textbf{which layer controls the steady-state flux}.
The fractional resistance of Layer~1 is $\eta = \kres/(1+\kres)$.
When $\kres \gg 1$, Layer~1 dominates (high-resistance barrier);
when $\kres \ll 1$, Layer~2 dominates.

\begin{definition}[Diffusion speed ratio]
\label{def:rho}
The \emph{diffusion speed ratio} is
\begin{equation}
  \label{eq:rho-def}
  \rspd \coloneqq \frac{\Ds{1}\,\Ls{2}}{\Ds{2}\,\Ls{1}}.
\end{equation}
\end{definition}

$\rspd$ compares the diffusive transit speeds of the two layers:
since the single-layer timelag scales as $\Ls{i}^2/(6\Ds{i})$,
$\rspd \gg 1$ means Layer~1 is ``fast'' (short transit time relative to its
thickness), while $\rspd \ll 1$ means Layer~1 is the diffusion bottleneck.

\begin{remark}
The two ratios encode all six parameters.  In particular,
$\kres = \Kpart/\rspd$ where $\Kpart = \Ss{2}/\Ss{1}$ is the partition
coefficient.  Thus any pair among $\{\kres, \rspd, \Kpart\}$ suffices
to characterise the bilayer, together with a single reference scale
(e.g., $\Ls{1}^2/\Ds{1}$ for time).
\end{remark}

\subsubsection{Dimensionless Timelag}

The composite timelag from \cref{eq:tlag-composite-raw} takes a transparent
form when non-dimensionalised as $\hat{\tau} = \tlag\,\Ds{1}/\Ls{1}^2$:
\begin{equation}
  \label{eq:tlag-dimensionless}
  \boxed{
  \hat{\tau}
  = \frac{\tfrac{1}{6} + \tfrac{\rspd^2}{2}
    + \tfrac{1}{\kres}\!\left(\tfrac{\rspd^2}{6}
    + \tfrac{1}{2}\right)}{1 + \tfrac{1}{\kres}}.}
\end{equation}
Each term in the numerator has a direct physical interpretation:
\begin{center}
\begin{tabular}{@{}ll@{}}
\toprule
\textbf{Term} & \textbf{Physical origin} \\
\midrule
$\tfrac{1}{6}$ & Layer~1 self-timelag
  ($\propto \Ls{1}^2/\Ds{1}$) \\[3pt]
$\tfrac{\rspd^2}{2}$ & Cross-loading: Layer~2 filled through Layer~1 \\[3pt]
$\tfrac{1}{\kres}\cdot\tfrac{\rspd^2}{6}$
  & Layer~2 self-timelag ($\propto \Ls{2}^2/\Ds{2}$) \\[3pt]
$\tfrac{1}{\kres}\cdot\tfrac{1}{2}$
  & Cross-loading: Layer~1 filled through Layer~2 \\
\bottomrule
\end{tabular}
\end{center}
The denominator $1 + 1/\kres$ is the total dimensionless resistance.
The cross-loading terms are the signature of the bilayer: they couple
information about \emph{both} layers into the timelag, even when one
layer dominates the self-timelag contribution.

\subsubsection{Regime Analysis}

Whether the extraction is well- or ill-conditioned depends on whether
\emph{both} $\tlag$ and $\Jss$ retain sensitivity to \emph{both} layers.
\Cref{tab:regimes} summarises the key regimes.

\begin{table}[ht]
\centering
\caption{Extraction conditioning as a function of the dimensionless ratios
$\kres$ and $\rspd$.}
\label{tab:regimes}
\begin{tabular}{@{}lccccl@{}}
\toprule
\textbf{Regime} & $\kres$ & $\rspd$
  & \textbf{$\tlag$ sees both?}
  & \textbf{$\Jss$ sees both?}
  & \textbf{Extraction} \\
\midrule
Both layers visible & $\sim$0.1--10 & $\sim$0.1--10
  & Yes & Yes & Well-conditioned \\[3pt]
Thin, fast Layer~1 & $\ll 1$ & $\gg 1$
  & No & No & Poor \\[3pt]
Thin barrier (Layer~1) & $\gg 1$ & any
  & Yes (cross-loading) & Yes (dominates) & Well-conditioned \\[3pt]
Identical materials & $=1$ & $=1$, $\Kpart=1$
  & Degenerate & Degenerate & Impossible \\[3pt]
One very slow layer & $\sim 1$ & $\ll 1$ or $\gg 1$
  & Partially & Yes & Marginal \\
\bottomrule
\end{tabular}
\end{table}

\paragraph{Why the thin, fast layer case fails.}
When $\kres \ll 1$ and $\rspd \gg 1$, Layer~1 is thin and diffuses rapidly.
Its self-timelag $\propto \Ls{1}^2/\Ds{1}$ is negligible, and its resistance
$R_1 = \Ls{1}/\Phi_1$ is a small fraction of the total.  Both the timelag
and the flux are dominated by Layer~2; changing $\Ds{1}$ or $\Ss{1}$ barely
alters either observable.  Physically, Layer~1 equilibrates almost
instantaneously and acts as a transparent window.

\paragraph{Why the fast/slow layer case is marginal.}
When $\rspd \gg 1$ with $\kres \sim 1$ (comparable resistances but very
different diffusivities), the timelag reduces to
$\hat{\tau} \approx \rspd^2/2 + (1/\kres)\rspd^2/6$, which depends on
$\rspd^2 = (\Ds{1}/\Ds{2})^2\cdot(\Ls{2}/\Ls{1})^2$.  The fast layer's
diffusivity $\Ds{1}$ cancels in the leading term after dividing by $d_0$:
\begin{equation}
  \tlag \approx \frac{\Ls{2}^2}{6\Ds{2}}
  \quad\text{when } \rspd \gg 1,\; \kres \sim 1.
\end{equation}
In other words, the timelag is blind to $\Ds{1}$ because
Layer~1 equilibrates instantaneously regardless of whether $\Ds{1}$ is
$10\times$ or $1000\times$ larger than $\Ds{2}$.  The flux $\Jss$ still
depends on $\Phi_1 = \Ds{1}\Ss{1}$ through the series resistance, providing
partial but not full information.

\subsubsection{The Thin Barrier Coating Case ($\kres \gg 1$)}
\label{sec:barrier-coating}

A particularly important practical scenario is a thin, low-permeability coating
(Layer~1) on a thick, more permeable substrate (Layer~2).  Here
$\Ls{1} \ll \Ls{2}$ but $\Phi_1 \ll \Phi_2$, so $\kres \gg 1$.
Despite the small thickness, the extraction is \textbf{well-conditioned}:

\paragraph{Flux sensitivity.}
Since $\kres \gg 1$, the barrier dominates the total resistance:
\begin{equation}
  \Jss \approx \frac{\sqrt{P}\,\Phi_1}{\Ls{1}}
  \qquad (\kres \gg 1).
\end{equation}
The flux is almost entirely determined by the coating's permeability $\Phi_1$,
providing high sensitivity to $\Ds{1}\Ss{1}$.

\paragraph{Timelag sensitivity.}
The composite timelag in this limit is (restoring dimensions):
\begin{equation}
  \label{eq:tlag-barrier-limit}
  \tlag \approx \frac{\Ls{1}^2}{6\Ds{1}}
  + \frac{\Kpart\,\Ls{1}\,\Ls{2}}{2\Ds{1}}
  + \frac{\Kpart\,\Ls{2}^2}{6\Ds{2}},
  \qquad (\kres \gg 1).
\end{equation}
The middle term---the \emph{cross-loading} contribution---is the key: it
represents the time required to load the thick substrate \emph{through} the
barrier bottleneck.  Even though $\Ls{1}$ is small, this term scales as
$\Ls{1}\Ls{2}/\Ds{1}$, and because $\Ls{2}$ is large and $\Ds{1}$ is small
(barrier property), it can dominate the timelag.
This ensures that $\Ds{1}$ remains visible in the timelag measurement
despite the barrier being thin.

\paragraph{Physical picture.}
In a barrier-coated membrane, hydrogen must pass through the thin but
impermeable coating to reach the thick substrate.  The substrate acts as a
large reservoir that must be ``filled'' through the coating bottleneck before
steady state is reached.  The filling time depends explicitly on $\Ds{1}$
(the coating diffusivity), keeping it experimentally accessible.

\paragraph{Exception.}
If $\kres \gg 1$ and $\rspd \gg 1$ \emph{simultaneously}, the layer has
high resistance \emph{because of very low solubility} ($\Ss{1} \ll \Ss{2}$)
rather than low diffusivity.  In this unusual case, the barrier diffuses
quickly but dissolves very little hydrogen; the cross-loading term in
\cref{eq:tlag-barrier-limit} becomes small because $\Kpart = \Ss{2}/\Ss{1}$
and $\Ds{1}$ are both large, and their ratio $\Kpart/\Ds{1}$ may not be.
This is atypical for practical barrier materials, which generally have
\emph{both} low $\Ds{i}$ and low $\Ss{i}$.

\subsubsection{Practical Diagnostic}

Before attempting extraction, a simple check can assess feasibility:

\begin{enumerate}
  \item \textbf{Compare timelags.}
    If a single-layer measurement of the substrate (Layer~2) is available,
    compute $\tau_2^* = (\Ls{2}^*)^2/(6\Ds{2})$.  If the bilayer timelag
    $\tlag$ differs significantly from $\tau_2^*\cdot(\Ls{2}/\Ls{2}^*)^2$
    (adjusting for thickness), then Layer~1 contributes measurably and
    extraction is feasible.
  \item \textbf{Estimate $\kres$.}
    Using literature values or order-of-magnitude estimates for $\Phi_i$,
    compute $\kres = (\Ls{1}/\Phi_1)/(\Ls{2}/\Phi_2)$.
    Values in the range $10^{-2} < \kres < 10^2$ (with $\kres$ not too close
    to~1 if $\Kpart \approx 1$) indicate a well-conditioned regime.
  \item \textbf{Estimate $\rspd$.}
    Compute $\rspd = \Ds{1}\Ls{2}/(\Ds{2}\Ls{1})$.
    The case $\kres \gg 1$ with $\rspd \lesssim 1$ (true barrier: low $D$,
    low $S$) is the most favourable for extraction.
\end{enumerate}

\begin{proposition}[Condition for well-posed extraction, refined]
\label{prop:well-posed}
The extraction of individual parameters is well-conditioned when either:
\begin{enumerate}
  \item Both layers contribute comparably to the total resistance
        ($\kres \sim 1$) and to the timelag ($\rspd \sim 1$), with the
        layers being materially distinct ($\Kpart \neq 1$ or $\rspd \neq 1$);
        or
  \item One layer dominates the resistance ($\kres \gg 1$ or $\kres \ll 1$)
        \textbf{and} the cross-loading term in the timelag remains
        significant---as is the case for thin barrier coatings where both
        $\Ds{i}$ and $\Ss{i}$ of the barrier are small.
\end{enumerate}
The extraction degenerates when both $\tlag$ and $\Jss$ lose sensitivity to
a layer simultaneously, which occurs for thin, fast, highly permeable layers
($\kres \ll 1$, $\rspd \gg 1$) or for identical materials
($\kres = \rspd = 1$, $\Kpart = 1$).
\end{proposition}

\subsection{Summary of Methods}
\label{sec:method-summary}

\begin{table}[ht]
\centering
\caption{Extraction methods for individual layer transport parameters.}
\label{tab:methods}
\begin{tabular}{@{}llcl@{}}
\toprule
\textbf{Method} & \textbf{Required data}
  & \textbf{Unknowns} & \textbf{Notes} \\
\midrule
1. Single-layer cal.\ + bilayer
  & $(\tau_i^*, J_i^*)$ + $(\tlag, \Jss)$
  & 4 & Most robust \\
2. Two bilayer samples
  & $(\tlag^{(k)}, \Jss^{(k)})$, $k=1,2$
  & 4 & No single-layer needed \\
3. Full transient fit
  & $\Cum(t)$ (full curve)
  & 4 & Single experiment \\
A. Known $\Ss{i}$: $\tlag + \Jss$
  & $\tlag$, $\Jss$
  & 2 & Simplest \\
B. Known $\Ss{i}$: $\tlag + \theta_{\mathrm{up}}$
  & $\tlag$, $\theta_{\mathrm{up}}$
  & 2 & No flux cal.\ needed \\
\bottomrule
\end{tabular}
\end{table}

\subsection{Practical Protocol}
\label{sec:protocol}

Based on the preceding analysis, we recommend the following experimental
procedure for extracting individual transport parameters from a bilayer
membrane when solubilities are \emph{not} known a priori.

\begin{enumerate}
  \item \textbf{Sample preparation.}  Fabricate (or obtain):
        \begin{itemize}
          \item A bilayer membrane with known $\Ls{1}$, $\Ls{2}$ (verified by
                cross-sectional microscopy if possible).
          \item \emph{Either} a single-layer sample of one constituent material
                (Method~1) \emph{or} a second bilayer sample with a different
                thickness ratio (Method~2).
        \end{itemize}

  \item \textbf{Single-layer calibration} (if using Method~1).
        Conduct a gas-driven permeation experiment on the single-layer sample:
        \begin{itemize}
          \item Extract $\tau_i^*$ and $\Jss^*$ from the $\Cum(t)$ asymptote.
          \item Compute $\Ds{i}$ and $\Ss{i}$ via \cref{eq:single-cal-D,eq:single-cal-Phi}.
        \end{itemize}

  \item \textbf{Bilayer permeation experiment(s).}
        \begin{itemize}
          \item Apply a known upstream hydrogen pressure~$P$.
          \item Record $\Cum(t)$ until well into steady state.
          \item Extract $\tlag$ (intercept) and $\Jss$ (slope) from the linear
                asymptote.
          \item Optionally, record the upstream pressure decay for
                $\theta_{\mathrm{up}}$ measurement.
        \end{itemize}

  \item \textbf{Parameter extraction.}
        \begin{itemize}
          \item \textbf{Method~1:} Use the known $\Ds{i}$, $\Ss{i}$ together
                with bilayer $\tlag$ and $\Jss$ to solve for $\Ds{j}$, $\Ss{j}$
                (\cref{sec:method1}).
          \item \textbf{Method~2:} Use both bilayer datasets to first determine
                $\Phi_1$, $\Phi_2$ from the flux system, then $\Ds{1}$, $\Ds{2}$
                from the two timelags (\cref{sec:method2}).
          \item \textbf{Method~3:} Fit the full transient $\Cum(t)$ to the
                analytical model (\cref{sec:method3}).
        \end{itemize}

  \item \textbf{Validation.}  Using the extracted parameters:
        \begin{itemize}
          \item Reconstruct the full transient $\Cum(t)$ curve using a FESTIM
                simulation or numerical inversion of the Laplace-domain
                solution.
          \item Compare the reconstructed curve with the measured data.
          \item Good agreement validates both the extracted parameters and the
                bilayer diffusion model; systematic discrepancies may indicate
                interface resistance, concentration-dependent diffusivity, or
                trapping effects.
        \end{itemize}
\end{enumerate}
